%%==========================================================================
\section{Introduction}
\label{sec:intro}
%%==========================================================================

The airline planning process is traditionally decomposed into a sequence of
sub-problems: timetable design, fleet assignment, crew rostering, tail
assignment, and recovery planning~\citep{barnhart2003applications,
gopalan1998mathematical}.  The \emph{tail assignment problem} (TAP)—assigning
a specific aircraft (identified by its tail number) to each flight leg—is the
penultimate planning step and is tightly coupled with aircraft maintenance
scheduling.  Poor tail assignment decisions directly inflate operational costs
through unnecessary ferry (repositioning) flights, suboptimal maintenance
routing, and safety-margin violations.

\citet{khaled2018compact} proposed a compact 0--1 linear programming
formulation for the TAP, featuring only $O(n_f n_p)$ binary variables (where
$n_f$ is the number of flights and $n_p$ the fleet size).  This is a
significant reduction over the classic multi-commodity flow model of
\citet{levin1971scheduling}, which requires $O(n_f^2 n_p)$ variables.
Computational experiments reported in that paper show that instances with up
to 1\,500 flight legs and 40 aircraft can be solved to certified optimality
within one hour on a standard workstation.

\subsection*{Our contributions}

The present paper makes four contributions that check, correct, and extend
those results:

\begin{enumerate}
\item \textbf{Basic feasibility correction (C2--C4 / constraints (3)--(6)).}
  We show that the strict-time balance definition used in the compact basic
  model admits a simultaneous-departure corner case: one aircraft can be
  assigned to two flights departing at the same airport and same time without
  violating the balance inequalities.  We formalize this counter-example and
  enforce physical executability by adding explicit pairwise non-overlap
  constraints.

\item \textbf{Constraint~(13) correction.}  We show that the cumulative-flight-
  time maintenance constraint (Constraint~13 in \citet{khaled2018compact})
  contains a logical flaw: the big-$M$ linearisation term $M(2 - y_{jd} -
  y_{jd'})$ allows the constraint to be vacuously satisfied even when the pair
  $(d, d')$ constitutes two consecutive maintenance events without an
  intervening check.  We provide a corrected formulation using two split
  endpoint constraints, prove its validity, and show that the corrected
  formulation yields strictly tighter LP relaxations on high-density instances.

\item \textbf{Full $\{A, B, C, D\}$ maintenance hierarchy.}  The model of
  \citet{khaled2018compact} handles type-A checks (cumulative flight hours).
  We extend the formulation to the full four-level hierarchy mandated by
  aviation regulators (FAA/EASA): type-A and~B checks are flight-hour
  triggered; type-C and~D checks are calendar-day triggered.  A heavier check
  subsumes all lighter checks (hierarchical counter resets).  We further add
  a constraint accounting for pre-horizon accumulated maintenance hours/days
  per aircraft.

\item \textbf{Heuristic benchmark.}  We implement and evaluate a family of
  constructive heuristics---a greedy/insertion baseline with repair and
  bounded local-search refinements---together with two space-time network
  search heuristics: a Dijkstra-based method (adapted from
  \citet{gronkvist2005tail}) and an ant-colony-optimization (ACO) variant.
  All heuristics respect the full maintenance hierarchy.  We report optimality
  gaps against the MILP on the same instance sets; the bounded local search
  improves the greedy seed by a mean of 12.7\% across the 25-instance suite,
  and the Dijkstra heuristic outperforms the ACO alternative in coverage on
  the space-time network.
\end{enumerate}

\subsection*{Paper organisation}

\Cref{sec:literature} reviews related work.  \Cref{sec:basic} restates the
basic TAP model and proves a C2--C4 corner-case correction.  \Cref{sec:maintenance}
presents the maintenance-extended model, the Constraint~(13) correction, and
the induced interpretation updates for C8--C14 under the corrected basic model.
\Cref{sec:hierarchy} describes the
full $\{A,B,C,D\}$ extension.  \Cref{sec:heuristics} analyses the heuristic
algorithms.  \Cref{sec:computational} reports all computational results.
\Cref{sec:conclusion} concludes.
